\begin{musica}{História de Portugal}{Tuna de Engenharia da Universidade do Porto}{}
\begin{multicols}{2}
\ac{Am}{Vou} contar aos igno\ac{E}{ran}tes
Os factos mais importantes
De que rezam a nossa his\ac{Am}{tó}ria;
\ac{Dm}{E} p’ra verem que não \ac{Am}{min}to
Ainda não toquei no \ac{E}{tin}to
P’ra não estragar a me\ac{Am}{mó}ria.

Se a memória não me falha
Foi numa grande batalha,
Lá no Estádio Nacional;
Que milhares de homens à bruta
Correram os filhos da luta
E lixaram Portugal.

Pedro Álvares Cabral
Foi parar ao hospital
Pelas escadinhas da bica,
Por um tal Egas Moniz
Lhe ter partido o nariz
Num Sporting – Benfica.

Dona Teresa queria gozo,
Enfeitava a testa ao esposo
Em bebedeiras de absinto;
Numa ruela d’Alfama
Dormiu com Vasco da Gama
Deu à luz D. João V.

Vieram então os reis de Espanha,
Que eram de primeira apanha
E falsos como os cães;
Puseram toda a cagança
Num bom ponta de lança
Que foi Fernão Magalhães.

Quando a construção civil
Foi importada do Brasil
Reinava D. Maria,
Que pela louça sanitária
Ficou sempre extraordinária
Rainha Maria Pia.

\ac{Am}{Mui}to mais havia a\ac{E}{in}da
Desta história tão linda
P’ra contar a toda a \ac{Am}{gen}te;
\ac{Dm}{Mas} ela só conti\ac{Am}{nua}
Quando o povo vier p’rá \ac{E}{rua}
E me fizer Presi\ac{Am}{den}te.
\ac{Dm}{Mas} ela só conti\ac{Am}{nua}
Quando o povo vier p’rá \ac{E}{rua}
E me fizer ENGE\ac{Am}{NHEI}RO! \ac{E Am}{}_ __
\end{multicols}

\posicoes{Dm: xx0231, E: 022100, Am: x02210}
\end{musica}
